\documentclass[10pt]{article}
\usepackage[a4paper,margin=1.6cm,top=1.4cm,bottom=1.5cm]{geometry}
\usepackage[T1]{fontenc}
\usepackage{lmodern}
\renewcommand{\familydefault}{\sfdefault}
\usepackage{xcolor,booktabs,tabularx,array,enumitem,titlesec,tcolorbox,hyperref}
\definecolor{brown}{HTML}{4B3527}
\definecolor{green}{HTML}{547965}
\definecolor{terra}{HTML}{AE664D}
\definecolor{sand}{HTML}{EFE9DF}
\hypersetup{colorlinks,urlcolor=green,linkcolor=green}
\titleformat{\section}{\large\bfseries\color{brown}}{}{0pt}{}
\titlespacing*{\section}{0pt}{6pt}{3pt}
\setlist[itemize]{leftmargin=1.2em,itemsep=1pt,topsep=2pt}
\setlength{\parindent}{0pt}
\setlength{\parskip}{2.5pt}
\newcommand{\rec}[2]{\par\vspace{3pt}{\color{brown}\bfseries\large #1.\ #2}\par\vspace{1pt}}
\newcommand{\conf}[1]{\textcolor{green}{\textbf{Confidence: #1}}}
\newcommand{\gate}[1]{\textbf{Go / no-go:} #1}
\newcommand{\act}[1]{\textbf{Do this:} #1}

\begin{document}

{\color{brown}\bfseries\LARGE Stories Coffee --- Recommendations}\\[2pt]
{\color{gray}Based on 51,177 item lines and 33,887 baskets across Garden, Harbour and Uptown, 1 Jan -- 30 Jun 2026 \textbullet\ \$305,925 revenue \textbullet\ \$179,430 recorded gross profit \textbullet\ 17 October 2026}

\begin{tcolorbox}[colback=sand,colframe=terra,boxrule=0.6pt,arc=2pt,left=6pt,right=6pt,top=4pt,bottom=4pt]
\textbf{Read this first.} The file looks simulated. Transaction IDs start with \texttt{SYN-}. Harbour's prices are exactly Garden's $\times 1.05$ and Uptown's $\times 0.96$. Prices and costs never change in 181 days. Hourly demand is pure random noise around a fixed weekly schedule. The March 2026 war, when Lebanon's restaurant syndicate reported falls of over 85\%, does not appear at all. Every recommendation below is therefore a \emph{tested pilot}, not a rollout. Reconcile the till with card settlements and cash deposits before committing capital.
\end{tcolorbox}

\section*{At a glance}
\renewcommand{\arraystretch}{1.12}
\begin{tabularx}{\linewidth}{@{}c>{\raggedright\arraybackslash}X>{\raggedright\arraybackslash}p{4.4cm}l@{}}
\toprule
\textbf{\#} & \textbf{Recommendation} & \textbf{Size of the prize / signal} & \textbf{Confidence}\\
\midrule
1 & Audit Garden's costs before changing suppliers & \$3,981 -- \$15,923 GP per 6 months & High (gap), cause open\\
2 & Staff, prep and seat by the clock, not the weather & Template within 1.2\% of best possible forecast & High\\
3 & Move the 10\% discount out of peak hours into quiet ones & Same budget; mornings $-37\%$, weekends $-29\%$ & Medium (test)\\
4 & Stop 8-day blanket 15\% campaigns; time them to slow periods & No measurable daily GP lift & Medium--high\\
5 & Run a no-discount evening attachment pilot: Cold Coffee + Bakery & Evening = 44.6\% of GP & Medium (test)\\
6 & Test a uniform +5\% price at one branch for four weeks & $\approx$\$7.3k GP per 6 months if demand holds & Medium (test)\\
7 & Test later opening at the weakest early hours & Median \$656 per 6 months, pays in 75\% of draws & Low--medium\\
\bottomrule
\end{tabularx}

\rec{1}{Audit Garden's costing rule before touching procurement}
Garden records a \textbf{46.91\%} gross margin, against Harbour's \textbf{64.96\%}, on the same 45 products.
\begin{itemize}
  \item Garden's unit cost is $\approx 1.43\times$ Harbour's on every product, and it never changes. That looks like an allocation rule, not waste.
  \item Repricing Garden's quantities at Harbour's costs adds \textbf{\$15,923.47} in six-month GP. If only 25--100\% of that is real cash, the range is \textbf{\$3,981 -- \$15,923}. If it is purely an accounting allocation, the saving is \$0.
\end{itemize}
\act{Trace the largest dollar contributors (the top five hold 22.3\% of the gap; Iced Matcha and the Chicken Pesto Sandwich lead) from the POS cost master through recipes to supplier invoices.}\\
\gate{change suppliers or recipes only if invoices show an avoidable cash difference larger than the cost of acting. At 50\% recovery the saving survives up to 14.5\% volume loss, so protect quality.} \conf{high that the gap exists; the cause is unknown.}

\rec{2}{Run operations off an hour $\times$ weekday template, not the weather}
Across all 81 features we engineered (weather, dust, holidays, Ramadan, school breaks, customer history), only the clock survives controls.
\begin{itemize}
  \item A branch $\times$ hour $\times$ weekday template predicts unseen May--June hours with an error of 1.59 baskets. A perfect model would still miss by 1.57, so the template is within \textbf{1.2\%} of the limit. Adding weather changes nothing (1.594).
  \item 19:00--24:00 brings \textbf{44.2\%} of baskets and \textbf{44.6\%} of GP. Dine-in jumps from \textbf{46\% to 72\%} at 18:00, and cold-drink share rises to \textbf{70\%} after 19:00, in every month and branch.
  \item The busiest slot is Harbour on Tuesday at 21:00: 21.3 expected baskets in June, \textbf{plan for 27} (90th percentile). Tuesday runs about $1.5\times$ other weekdays.
\end{itemize}
\act{roster staff, ice, cold-brew prep, desserts and seating from the staffing template (in the project database). Log queues and stock-outs from 19:00 to 23:00 for four weeks.}\\
\gate{add evening labour only if logged lost orders $\times$ \$5.34 GP per basket exceed its cost.} \conf{high.}

\rec{3}{Spend the existing 10\% discount where demand is weak}
Today the discount ignores the clock. About \textbf{5.5\%} of items are discounted in the morning, afternoon and evening alike, and on every weekday (5.1--6.1\%).
\begin{itemize}
  \item Mornings bring \textbf{37\%} fewer baskets per hour and weekends \textbf{29\%} fewer per day. Yet over half of discounted items (52\%) are sold in the busy evening, to customers who would come anyway.
  \item The payday window (last 3 and first 3 days of the month) already brings +4.7\% daily baskets, so don't discount into it.
\end{itemize}
\act{for four weeks, at one branch, give the 10\% offer only before 11:00 and on Saturdays. Keep the other branches as controls.}\\
\gate{keep it only if GP per assigned visitor rises against control. Basket averages mislead, because bigger baskets are more likely to contain a discounted item by chance.} \conf{medium.}

\rec{4}{Retire blanket 8-day 15\% campaigns; rerun them as a timed experiment}
The two 15\% windows, 12--19 March and 12--19 May, ran at Harbour on Cold Coffee and Desserts across every weekday and at peak hours.
\begin{itemize}
  \item Against Garden and Uptown over the same dates, Harbour's daily GP during the campaigns was \emph{lower}, with wide intervals. Against Harbour's own no-discount baskets, the 15\% baskets earned \textbf{\$0.33} less GP.
\end{itemize}
\act{randomize 0\% / 10\% / 15\% offers by day at Harbour and log people who don't buy.}\\
\gate{scale only an arm that raises contribution per assigned visitor after discount cost.} \conf{medium--high that the current format doesn't pay.}

\rec{5}{Pilot evening attachment without a discount: Cold Coffee + Bakery}
The evening is a sit-down cold-drink and dessert occasion: Cold Coffee is 16.6\% of evening units, Matcha 17.5\% and Desserts 10.7\%. Dine-in and takeaway baskets earn the same GP (\$5.32 vs \$5.38), so the gain must come from more items per basket, not from seating.
\begin{itemize}
  \item Cold Coffee + Bakery repeats in both halves of the period, carries a 66\% recorded margin, and coffee input costs are falling. Matcha supply is being rationed and its prices set records, so don't stimulate demand for it before securing supply.
\end{itemize}
\act{a scripted ``add a pastry?'' prompt from 19:00 to 23:00 with no discount, run on alternating weeks.}\\
\gate{keep it if attached GP exceeds labour time and any service slowdown at the peak.} \conf{medium.}

\rec{6}{Test a uniform +5\% price rise at one branch}
At Stories' 58.7\% margin, a 5\% rise pays unless demand falls by more than 1.68\% for every 1\% of price. A 20,000-draw simulation puts the gain at \textbf{\$7,305} per six months (5th--95th percentile \$3,025 -- \$11,128). That certainty comes from the assumptions, not from data: prices never changed in the file.
\begin{itemize}
  \item Holding coffee prices back loses to a uniform rise in 88\% of draws. Eleven Harbour prices sit a few cents above a whole dollar, the worst place for a price.
\end{itemize}
\act{+5\% at one branch for four weeks against the other two, with price endings cleaned up.}\\
\gate{stop if baskets per day fall by more than the break-even drop of about 8\% over the test period.} \conf{medium.}

\rec{7}{Test later opening at the weakest early hours}
Uptown's 07:00 hour averages \$4.37 GP against an estimated \$6.53 in staff and generator cost. The generator tariff rose 74\% from January to April. Closing Uptown 07:00--09:00 and Garden 07:00 pays in 75\% of simulated scenarios, with a median of \$656 per six months.
\begin{itemize}
  \item Staffing assumptions drive this result far more than energy prices. Confirm the real early-shift roster first.
\end{itemize}
\act{open Uptown at 09:00 and Garden at 08:00 for four weeks.}\\
\gate{reverse it if those customers don't come back later in the day.} \conf{low--medium.}

\section*{What we would \emph{not} do --- and why}
\begin{itemize}
  \item \textbf{End-of-day discounts.} Late evening is Stories' \emph{peak}, not its slack. Discounting it gives margin away to customers who are already coming.
  \item \textbf{A seasonal ``summer iced'' push or a weather-data subscription.} Cold-drink share stays at 60.2--60.9\% as Beirut warms from 14.0\,\textdegree C to 25.5\,\textdegree C. No weather, dust or holiday effect survives controls.
  \item \textbf{Winter or holiday promotions.} January is the weakest month, but 99.2\% of the growth to June came from more baskets per day, not bigger baskets, and holidays show no measurable effect.
  \item \textbf{Bundle discounts based on how often items are bought together.} Pairs bought together don't show that a discount would add sales, and bigger baskets carry discounts by chance.
\end{itemize}

\vspace{2pt}
{\small\color{gray}Evidence: the reports feature\_loop, external\_factors, executive\_summary, outside\_in\_chains and monte\_carlo, plus the discount study in analysis/outputs; all tables are in reports/tables. Effects are associations measured in historical data; financial figures are scenarios before implementation costs, not forecasts.}

\end{document}
